\section{方法选择与结果检验}
\label{sec:ls-direct-method-comparison}

前面各节分别给出了求最小二乘解与最小范数解的直接法。面对具体问题时，要先回到第~\ref{chap:linear-solution-selection}章确定的求解目标，再看矩阵的形状与秩，最后用原始数据核验算出的结果。本节按这一顺序归纳，并在同一组数据上比较几条路线。

\paragraph{按目标选择}
求最小二乘解时，满列秩用QR，条件数适中、需要按行累计数据时也可用正规方程法；满行秩欠定方程的最小范数解用转置QR；秩亏或秩未知的一般情形用SVD。加权目标先白化再按普通情形求解，齐次目标取最小伸缩方向。方阵且相容的精确解问题则回到第~\ref{chap:direct-linear-solvers}章，按矩阵结构选择LU、Cholesky或分块LDL。

\paragraph{同一组数据上的比较}
沿用示例~\ref{ex:ls-line-least-squares}的直线拟合$\bm A=\left[\begin{smallmatrix}1&0\\1&1\\1&2\end{smallmatrix}\right]$、$\bm b=(1,2,2)^\top$。正规方程法（示例~\ref{ex:ls-normal-equations-line-fit}）形成$\bm G=\left[\begin{smallmatrix}3&3\\3&5\end{smallmatrix}\right]$，得到Cholesky因子$\bm L$与前代结果$\bm y=(5/\sqrt3,1/\sqrt2)^\top$；Givens路线（示例~\ref{ex:ls-givens-line-fit}）得到$\bm R=\bm L^\top$与$\bm c_1=\bm y$，Householder路线（示例~\ref{ex:ls-householder-line-fit}）得到的$\bm R$与$\bm c_1$只差各行的符号。三条路线最后回代的是同一个三角方程，给出同一个$\hat{\bm x}=(7/6,1/2)^\top$；区别在得到这个方程的方式，正规方程法要先形成$\bm A^\top\bm A$，正交方法不需要。

复制常数列后得到$\widetilde{\bm A}=\left[\begin{smallmatrix}1&0&1\\1&1&1\\1&2&1\end{smallmatrix}\right]$，列空间和最优残差都不变，但秩降为$2$，正规方程与QR回代都会遇到零主元。SVD求解略去$\sigma_3=0$对应的方向$(1,0,-1)^\top/\sqrt2$，得到最小范数解$\bm x_\star=(7/12,1/2,7/12)^\top$：最小二乘只约束$x_1+x_3=7/6$与$x_2=1/2$，最小范数再把$7/6$平分给两个相同的列。前一组比较的差别在计算方式，这里的差别则在选解目标。

\paragraph{用原始数据核验}
算出$\bm x$后，要代回原问题检查它是否符合目标，而不只检查中间的小系统。普通最小二乘解检查$\bm A^\top(\bm b-\bm A\bm x)=\bm0$，要求最小范数时再检查$\bm x\perp\ker\bm A$。在上例中，残差$(-\tfrac16,\tfrac13,-\tfrac16)^\top$不为零，但它与$\bm A$的两列都正交；$\widetilde{\bm A}\bm x_\star=\bm A\hat{\bm x}$，且$\bm x_\star$与核方向$(1,0,-1)^\top$正交。加权问题使用第~\ref{sec:ls-least-squares-verification}节的加权正交条件$\bm A^\top\bm M(\bm b-\bm A\bm x)=\bm0$，齐次问题检查单位长度以及是否为最小伸缩方向。最小二乘的残差本来就可以非零；这些量用于核对计算过程，各求解算法仍只返回$\bm x$。

计算中还可以用分解残差（如$\lVert\bm A-\bm Q_1\bm R\rVert$）和正交性（如$\lVert\bm Q_1^\top\bm Q_1-\bm I\rVert$）分别核对中间变换，但这些检查只是旁证，不属于算法本身，也不能代替对原问题的检验。

\paragraph{残差小不等于误差小}
即使上述检验都通过，也不保证$\bm x$的误差小：矩阵接近奇异时，很小的残差可能对应很大的解误差。本章只给出了算法层面的结论：正规方程使条件数平方，正交变换保持长度，SVD把小奇异值的放大暴露为$1/\sigma_r$。第~\ref{chap:solution-error-regularization}章用条件数与前向、后向误差系统地解释这些现象（第~\ref{sec:ls-numerical-sensitivity}节），并在直接法与迭代法之间作整体比较（第~\ref{sec:ls-method-selection}节）。
